\documentclass{article}

\usepackage{amsmath,amssymb}
\usepackage{xurl}
\usepackage[hidelinks]{hyperref}
\setlength{\emergencystretch}{2em}

% Shared commands for notes in docs/paper-gaps/.

\DeclareUnicodeCharacter{2080}{\ifmmode{}_{0}\else\textsubscript{0}\fi}
\DeclareUnicodeCharacter{2081}{\ifmmode{}_{1}\else\textsubscript{1}\fi}
\DeclareUnicodeCharacter{2082}{\ifmmode{}_{2}\else\textsubscript{2}\fi}
\DeclareUnicodeCharacter{2099}{\ifmmode{}_{n}\else\textsubscript{n}\fi}

\newcommand{\Lean}{\textsc{Lean}}
\ExplSyntaxOn
\NewDocumentCommand{\leanid}{m}{
  \ifmmode
    \textsf{\detokenize{#1}}
  \else
    \group_begin:
    \urlstyle{sf}
    \tl_set:Nn \l_tmpa_tl {#1}
    \tl_replace_all:Nnn \l_tmpa_tl {\_} {_}
    \exp_args:NV \path \l_tmpa_tl
    \group_end:
  \fi
}
\ExplSyntaxOff
\providecommand{\tr}{\operatorname{tr}}
\newcommand{\tnleanIssueBase}{https://github.com/LionSR/TNLean/issues/}
\newcommand{\tnleanPrBase}{https://github.com/LionSR/TNLean/pull/}
\newcommand{\tnleanDocBase}{https://sirui-lu.com/TNLean/docs/}
\newcommand{\ghissue}[1]{\href{\tnleanIssueBase#1}{GitHub issue \##1}}
\newcommand{\ghpr}[1]{\href{\tnleanPrBase#1}{PR \##1}}
\newcommand{\doclink}[1]{\href{\tnleanDocBase#1.html}{\path{#1}}}

% Dirac notation and the identity operator, as in blueprint/src/macros/common.tex.
% \providecommand keeps note-local definitions authoritative.
\usepackage{dsfont}
\providecommand{\ket}[1]{|#1\rangle}
\providecommand{\bra}[1]{\langle #1|}
\providecommand{\braket}[2]{\langle #1 | #2 \rangle}
\providecommand{\ketbra}[2]{|#1\rangle\!\langle #2|}
\providecommand{\Id}{\mathds{1}}
\providecommand{\one}{\mathds{1}}
\providecommand{\C}{\mathbb{C}}
\providecommand{\MN}[1]{M_{#1}(\C)}
\providecommand{\MD}{M_D(\C)}

% Verdict marker: \gapnote{<kind>}{<status>}. Kind is the policy
% classification (clarification, local-correction, scope-restriction,
% unfaithful, false-source, open-gap); status is open, wip (elimination
% underway), resolved, or historical. Machine-read by the site index;
% prints a verdict line.
\newcommand{\gapnote}[2]{\par\noindent\textbf{Verdict:} #1 (#2).\par}


\title{Scope of the Formal Site-Independent Construction Relative to PGVWC07}
\author{}
\date{2026-10-29}

\begin{document}
\maketitle
\gapnote{scope-restriction}{open}

\section{Source theorem}

Theorem~3 of Ref.~\cite{PerezGarcia2007Matrix} states that every
translation-invariant pure state on a finite periodic chain has an MPS
representation with site-independent matrices. Let $N>0$ be the ring length.
The source first applies successive Schmidt or singular value decompositions
to obtain a varying-bond open-boundary representation. It then embeds the
rectangular virtual spaces into a common square virtual space and averages the
padded chain over cyclic translations. If $D$ bounds the intermediate bond
dimensions, the resulting site-independent tensor has bond dimension
\begin{align}
    D_{\mathrm{site}}
     & = ND.
    \label{eq:pgvwc07_site_independent_scope_site_independent_bond_dimension}
\end{align}

\section{Formal boundary}

The structure \leanid{OBCChainTensor} records a supplied varying-bond
open-boundary representation. For a fixed positive length $N$, the theorem
\leanid{OBCChainTensor.pgvwc07_site_independent_matrices} proves that cyclic
averaging of the zero-padded representation preserves its coefficients when
those coefficients are translation invariant. Its existential form is
\leanid{OBCChainTensor.exists_pgvwc07_site_independent_tensor}.

These declarations prove the square-padding and cyclic-averaging part
with a supplied OBC representation. The raw arbitrary-state OBC existence
step is already available through
\leanid{OBCChainTensor.exists_coeff_eq_of_cutRank_le}, with the actual
physical cut ranks and the bound $d^{\lfloor N/2\rfloor}$.

\section{The composed construction now written}

The theorem \leanid{MPSTensor.exists_mpv_eq_of_cyclic_invariant}
composes these existing constructions. Given a nonzero physical vector
on $N>0$ sites, fixed exactly by every cyclic translation, it constructs
an OBC chain with bound $d^{\lfloor N/2\rfloor}$, derives translation
invariance of its coefficient function, and constructs site-independent
matrices of size $N d^{\lfloor N/2\rfloor}$. No OBC representation is a
premise. The Hilbert-space version is
\leanid{MPSTensor.exists_mpvState_eq_of_cyclic_invariant}.

The tensor still depends on $N$ and agrees with the original vector at
that length; no other-length equality is asserted. Exact vector
translation invariance is used, as in the source's cyclic averaging.
Invariance of a density matrix or a one-dimensional state line alone
allows a nontrivial cyclic eigenphase and does not imply this vector
identity. The exact representability criterion is
\leanid{MPSTensor.exists_mpvState_eq_iff_cyclic_invariant}.

\section{Verdict}

\textbf{Verdict: composed proof written; verification pending.}
The previously documented supplied-OBC boundary is removed in the new
source-facing composed theorem. The original scoped lemma is retained
as the useful finite cyclic-block calculation. The new declaration and
Blueprint entry remain pending hosted Lean verification and integration;
no local Lean invocation was run. The historical note stays open until
that acceptance is complete.

\bibliographystyle{alpha}
\bibliography{references}

\end{document}